
\subsubsection{Research Questions}

We structure our proof-of-concept validation around three technical questions:

\textbf{Q1 (Structural Contracts)}: Can decorator-based structural contracts detect shape and dtype mismatches at import time with zero false positives?

\textbf{Q2 (Metamorphic Contracts)}: Can metamorphic contracts detect softmax sum and dropout identity violations with $\geq 70$\% detection rate and $\leq 10s$ execution time?

\textbf{Q3 (Composition Contracts)}: Can cross-framework validation contracts detect dtype corruption, shape inconsistency, and numerical drift with $\geq 70$\% detection rate, <10\% false positive rate, and <5s validation time?

\subsubsection{Experimental Design}

\textbf{Q1: Structural Contracts (h-m1)}

\textbf{Protocol}: Controlled scenario validation with real tensor operations.

\textbf{Procedure}:
\begin{enumerate}
\item Implement \texttt{@validate\textbackslash \_structural} decorator for tensor shape and dtype validation
\item Apply decorator to model initialization
\item Test scenarios:
\begin{itemize}
\item Control: 3-channel input, model expects 3 channels (should pass)
\item Defect 1: Model expects 4 channels, dataset provides 3 (should detect)
\item Defect 2: Model expects float32, receives float16 (should detect)
\item Measure detection rate and execution time
\end{itemize}
\end{enumerate}

\textbf{Success Criterion}: 100\% detection on defects, 0\% false positives on control, <0.1s execution time.

\textbf{Q2: Metamorphic Contracts (h-m2)}

\textbf{Protocol}: Synthetic defect injection with controlled ground truth.

\textbf{Dataset}:
\begin{itemize}
\item Total: 50 scenarios
\item Control: 10 (valid operations, no violations)
\item Softmax violations: 20 (perturbed outputs, sum $\neq$ 1.0)
\item Dropout violations: 20 (identity broken in eval mode)
\item Seed: 42 (deterministic, reproducible)
\end{itemize}

\textbf{Baseline}: Plain PyTorch operations without metamorphic validation (expected 0\% detection).

\textbf{Proposed}: Metamorphic contracts with:
\begin{itemize}
\item Softmax sum validation: \texttt{torch.allclose(sum, 1.0, rtol=1e-5, atol=1e-7)}
\item Dropout identity validation: \texttt{torch.equal(dropout(x, eval=True), x)}
\end{itemize}

\textbf{Metrics}:
\begin{itemize}
\item Detection rate: \% of injected violations detected
\item False positive rate: \% of control scenarios flagged
\item Execution time: average validation time per scenario
\end{itemize}

\textbf{Success Criterion}: Detection rate $\geq 70$\%, false positive rate <5\%, execution time $\leq 10s$ total.

\textbf{Q3: Composition Contracts (h-c2)}

\textbf{Protocol}: Synthetic cross-framework defect injection.

\textbf{Dataset}:
\begin{itemize}
\item Total: 30 test cases (70/30 train-test split from 100 synthetic samples)
\item Defects: 24 (numerical drift: 10, shape inconsistency: 9, dtype corruption: 5)
\item Valid conversions: 6
\end{itemize}

\textbf{Defect Injection}:
\begin{itemize}
\item Dtype corruption: Force float16 output instead of float32
\item Shape mismatch: Truncate output dimension
\item Numerical drift: Add Gaussian noise ($\sigma =0.1)$
\end{itemize}

\textbf{Contracts}:
\begin{enumerate}
\item Dtype preservation: Validate output dtype matches source
\item Shape consistency: Validate output shape equivalence
\item Numerical tolerance: Validate numerical equivalence (rtol=1e-5, atol=1e-7)
\end{enumerate}

\textbf{Baseline}: End-to-end validation (run inference, compare outputs with np.allclose).

\textbf{Metrics}:
\begin{itemize}
\item Detection rate: \% of defects detected
\item False positive rate: \% of valid conversions flagged
\item Validation time: average time per model
\end{itemize}

\textbf{Success Criterion}: Detection rate $\geq 70$\%, false positive rate <10\%, validation time <5s.

\subsubsection{Limitations}

\textbf{Sample Size}: All experiments use small sample sizes (N=2 for structural, N=50 for metamorphic, N=30 for composition). These are proof-of-concept validations demonstrating technical feasibility, not statistical validation for production deployment.

\textbf{Synthetic Defects}: Metamorphic and composition experiments use programmatically injected defects rather than real-world bugs from production repositories. Synthetic defects enable controlled validation but may not represent the full distribution of failure modes in practice.

\textbf{Domain}: All experiments use PyTorch operations and computer vision patterns. Generalization to other frameworks (TensorFlow, JAX) and domains (NLP, RL) remains untested.

\textbf{Version Stability}: None of the proof-of-concept experiments test contracts across multiple library versions. Version stability is a critical requirement for production deployment but is deferred to future work.
